\heading{实验物理中的统计方法-参数估计的一般概念}





\begin{question}
考虑随机数 \(x\)，其均值和方差分别为 \(\mu\) 和
\(\sigma^2\)。假设样本空间由 \(n\) 次观测结果 \(x_1,x_2,\cdots,x_n\)
构成。本题的目的是证明样本均值

\[
\bar x=\frac1n\sum_{i=1}^{n}x_i
\]

为均值 \(\mu\) 的一致性估计量。

\begin{enumerate}
\item
  第一步要证明 Chebyshev 不等式，即只要 \(x\) 的方差存在，对任意
  \(a>0\)，下面式子成立：
\end{enumerate}

\[
P(|x-\mu|\ge a)\le \frac{\sigma^2}{a^2}.
\]

该式的证明需要用到方差的定义，

\[
\sigma^2=\int_{-\infty}^{\infty}(x-\mu)^2f(x)\,dx,
\]

其中 \(f(x)\) 为随机变量 \(x\)
的概率密度函数。利用如下事实：如果积分区域限定为
\(|x-\mu|\ge a\)，则积分 (5.3) 会变小，如果用 \(a^2\) 替换
\((x-\mu)^2\)，则积分会更小。

\begin{enumerate}[start=2]
\item
  利用 Chebyshev 不等式证明大数弱定理，即，对任意 \(\epsilon>0\)，
\end{enumerate}

\[
\lim_{n\to\infty}P\left(\left|\frac1n\sum_{i=1}^{n}x_i-\mu\right|\ge\epsilon\right)=0.
\]

这等价于，\(\bar x\) 是 \(\mu\) 的一致性估计量，只要 \(x\) 的方差存在。
\begin{solution}
设随机变量 \(x\) 的均值与方差分别为 \(\mu,\sigma^2\)，样本均值为

\[
\bar x=\frac1n\sum_{i=1}^n x_i.
\]

\begin{enumerate}
\item 证明 Chebyshev 不等式

由方差定义

\[
\sigma^2=E[(x-\mu)^2]=\int_{-\infty}^{\infty}(x-\mu)^2 f(x)\,dx.
\]

在区域 \(|x-\mu|\ge a\) 上，恒有 \((x-\mu)^2\ge a^2\)，于是

\[
\sigma^2 \ge \int_{|x-\mu|\ge a}(x-\mu)^2f(x)\,dx
\ge a^2\int_{|x-\mu|\ge a}f(x)\,dx.
\]

因此

\[
P(|x-\mu|\ge a)\le \frac{\sigma^2}{a^2}.
\]

这就是 Chebyshev 不等式。

\item 用 Chebyshev 证明大数弱定理

若 \(x_1,\dots,x_n\) 独立同分布，均值为 \(\mu\)，方差为 \(\sigma^2\)，则

\[
E[\bar x]=\mu,\qquad V[\bar x]=\frac{\sigma^2}{n}.
\]

对任意 \(\varepsilon>0\)，由 Chebyshev 不等式，

\[
P(|\bar x-\mu|\ge \varepsilon)
\le \frac{V[\bar x]}{\varepsilon^2}
=\frac{\sigma^2}{n\varepsilon^2}
\xrightarrow[n\to\infty]{}0.
\]

故 \(\bar x\) 依概率收敛到 \(\mu\)，即 \(\bar x\) 是 \(\mu\)
的一致估计量。
\end{enumerate}
\end{solution}
\end{question}



\begin{question}
考虑均值为 \(\mu\)，方差为 \(\sigma^2\) 的随机变量 \(x\)，并得到样本值为
\(x_1,x_2,\cdots,x_n\) 的样本空间。

\begin{enumerate}
\item
  假设均值 \(\mu\) 已经利用样本均值 \(\bar x\) 估计。证明样本方差
\end{enumerate}

\[
s^2=\frac{1}{n-1}\sum_{i=1}^{n}(x_i-\bar x)^2=\frac{n}{n-1}(\overline{x^2}-\bar x^2)
\]

为方差 \(\sigma^2\) 的无偏估计量。（利用
\(E[x_ix_j]=\mu^2\ (i\ne j)\)，\(E[x_i^2]=\mu^2+\sigma^2\ (i=1,2,\ldots,n)\)。）

\begin{enumerate}[start=2]
\item
  假设均值 \(\mu\) 已知。证明 \(\sigma^2\) 的无偏估计量为
\end{enumerate}

\[
S^2=\frac1n\sum_{i=1}^{n}(x_i-\mu)^2=\overline{x^2}-\mu^2.
\]
\begin{solution}
样本方差定义为

\[
s^2=\frac1{n-1}\sum_{i=1}^n (x_i-\bar x)^2.
\]

先用恒等式

\[
\sum_{i=1}^n (x_i-\bar x)^2=\sum_{i=1}^n (x_i-\mu)^2-n(\bar x-\mu)^2.
\]

两边取期望：

\[
E\left[\sum_{i=1}^n (x_i-\bar x)^2\right]
= n\sigma^2 - nV(\bar x)
= n\sigma^2 - n\frac{\sigma^2}{n}
= (n-1)\sigma^2.
\]

因此

\[
E[s^2]=\sigma^2.
\]

所以 \(s^2\) 是无偏估计量。

\begin{enumerate}[start=2]
\item 若均值 \(\mu\) 已知，证明

\[
S^2=\frac1n\sum_{i=1}^n (x_i-\mu)^2
\]

是 \(\sigma^2\) 的无偏估计。

直接取期望：

\[
E[S^2]=\frac1n\sum_{i=1}^n E[(x_i-\mu)^2]=\frac1n\sum_{i=1}^n \sigma^2=\sigma^2.
\]

因此 \(S^2\) 无偏。
\end{enumerate}
\end{solution}
\end{question}



\begin{question}
\begin{enumerate}
\item
  证明样本方差 \(s^2\) (5.5) 的方差为
\end{enumerate}

\[
V[s^2]=E[s^4]-(E[s^2])^2=\frac1n\left(\mu_4-\frac{n-3}{n-1}\mu_2^2\right),
\]

其中 \(\mu_k=E[(x-\mu)^k]\) 为 \(x\) 的 \(k\) 阶中心矩。为此，需要先证明
\(s^2\) 可以写为

\[
s^2=\frac{1}{n-1}\sum_{i=1}^{n}x_i^2-\frac{1}{n(n-1)}\sum_{i,j=1}^{n}x_ix_j.
\]

然后证明 \(s^4\) 的期待值为

\[
\begin{split}
E[s^4]=&\frac{1}{(n-1)^2}\sum_{i,j=1}^{n}E[x_i^2x_j^2]
-\frac{2}{n(n-1)^2}\sum_{i,j,k=1}^{n}E[x_ix_jx_k^2]\\
&+\frac{1}{n^2(n-1)^2}\sum_{i,j,k,l=1}^{n}E[x_ix_jx_kx_l].
\end{split}
\]

计算每个求和给出多少项代数矩 \(\mu'_4\) 或
\(\mu_2'^2\)。注意其余所有项都含 \(\mu\) 的一次项或高次项。令
\(\mu=0\)，将结果表示为中心矩 \(\mu_2\) 和 \(\mu_4\)
的函数。从中减掉本章第 2 题得到的 \((E[s^2])^2\) 即可得到结果。

\begin{enumerate}[start=2]
\item
  如果 \(x\) 服从高斯分布，计算 \(s^2\) 的方差。利用高斯分布的 4
  阶中心矩为 \(\mu_4=3\sigma^4\)。
\end{enumerate}
\begin{solution}
\[
V[s^2]=E[s^4]-(E[s^2])^2=
\frac1n\left(\mu_4-\frac{n-3}{n-1}\mu_2^2\right),
\]

其中 \(\mu_k=E[(x-\mu)^k]\) 是中心矩。

把随机变量平移到均值为 0 的情形（不影响方差与中心矩），写成

\[
s^2=\frac1{n-1}\sum_{i=1}^n x_i^2-
\frac1{n(n-1)}\sum_{i,j=1}^n x_ix_j.
\]

接着平方并对所有指标重合方式分类求期望：

\begin{itemize}
\item
  四个指标全不相同；
\item
  有两个指标重合；
\item
  三个重合；
\item
  两对重合；
\item
  四个全重合。
\end{itemize}

由于独立同分布且 \(E[x]=0\)，许多项会消失，只剩下由 \(\mu_2=E[x^2]\) 与
\(\mu_4=E[x^4]\) 组成的项。整理后可得

\[
E[s^4]=\frac1n\mu_4+
\frac{n^2-2n+3}{n(n-1)}\mu_2^2,
\]

再减去

\[
(E[s^2])^2=\mu_2^2,
\]

便得到题目给出的结果：

\[
V[s^2]=\frac1n\left(\mu_4-\frac{n-3}{n-1}\mu_2^2\right).
\]

\begin{enumerate}[start=2]
\item 若 \(x\) 服从高斯分布

高斯分布满足

\[
\mu_2=\sigma^2,\qquad \mu_4=3\sigma^4.
\]

代回上式：

\[
V[s^2]
=\frac1n\left(3\sigma^4-\frac{n-3}{n-1}\sigma^4\right)
=\frac{2\sigma^4}{n-1}.
\]
\end{enumerate}
\end{solution}
\end{question}

